\documentclass[../../../include/open-logic-section]{subfiles}

\begin{document}

\olfileid{sth}{z}{power}
\olsection{幂集}

下面我们继续引入另一条公理：

\begin{axiom}[幂集公理]
对于任何集合 $A$，集合 $\Pow{A} = \Setabs{x}{x \subseteq A}$ 存在。
\[
	\forall A \exists P \forall x(x \in P \liff (\forall z \in x)z \in A)
\]
\end{axiom}

对此的论证相当直接。假设 $A$ 在阶段 $S$ 形成。那么 $A$ 的所有成员在 $S$ 之前就已经可用（根据 \stagesacc）。因此，与分离公理的论证同理，$A$ 的每个子集在阶段 $S$ 就已形成。所以在 $S$ 之后的任何阶段，它们都是可用的，从而能构成一个单一集合。并且我们知道存在某个这样的阶段，因为 $S$ 不是最后一个阶段（根据 \stagessucc）。所以 $\Pow{A}$ 存在。

幂集公理有一个很好的推论：

\begin{prop}\label{thm:Products}
给定任意集合 $A, B$，它们的笛卡尔积 $A \times B$ 存在。
\end{prop}

\begin{proof}
根据幂集公理和 \olref[sth][z][pairs]{prop:pairsconsequences}，集合 $\Pow{\Pow{A \cup B}}$ 存在。因此，由分离公理，下述集合存在：
\[
	C = \Setabs{z \in \Pow{\Pow{A \cup B}}}{(\exists x \in A)(\exists y \in B) z = \tuple{x, y}}.
\]
对于任意 $x \in A$ 和 $y \in B$，根据 \olref[sth][z][pairs]{prop:pairsconsequences}，集合 $\tuple{x, y}$ 存在。此外，由于 $x, y \in A \cup B$，我们有 $\{x\}, \{x, y\} \in \Pow{A \cup B}$，并且 $\tuple{x,y} \in \Pow{\Pow{A \cup B}}$。所以 $A \times B = C$。
\end{proof}

在这个证明中，幂集公理与分离公理配合使用。这并不意外。没有分离公理，幂集公理就不会是一条非常\emph{强有力}的原理。毕竟，分离公理告诉我们一个集合的哪些子集存在，从而决定了每个幂集有多“肥”。

\begin{prob}
证明：对于任意集合 $A, B$，(i) 所有定义域为 $A$、值域为 $B$ 的关系构成的集合存在；(ii) 所有从 $A$ 到 $B$ 的函数构成的集合存在。
\end{prob}

\begin{prob}
设 $A$ 为集合，$\sim$ 为 $A$ 上的等价关系。证明 $A$ 在 $\sim$ 下的等价类所构成的集合，即 $\equivclass{A}{\sim}$，存在。
\end{prob}

\end{document}